% !TEX root = ../../main.tex
% ============================================================================
% Chapter: Clutching uniqueness of the Hodge Euler class on the socle
%
% Platonic-ideal inscription of the AP225 uniqueness statement.
% This chapter SUPPLEMENTS prop:scalar-obstruction-hodge-euler
% (which is already unconditional via Arakelov-Faltings + BGS + GRR)
% with a separate boundary-restriction uniqueness statement, stated at the
% honest scope: uniqueness holds on the SOCLE of the tautological ring,
% not on the full ring at degree g > 1.
%
% Guiding principle: strengthening only. Every claim is scoped to the
% precise invariant (socle projection) that Graber-Vakil type theorems
% actually control. No bare kappa (AP113). All labels unique (AP124/125).
% No AI slop (AP29).
% ============================================================================

\providecommand{\kscal}{\kappa}
\providecommand{\kMum}{\kappa_{1}}
\providecommand{\cA}{\mathcal{A}}

\chapter{Clutching uniqueness on the socle of the tautological ring}
\label{chap:clutching-uniqueness-socle}
\index{clutching!uniqueness on the socle|textbf}
\index{tautological ring!socle projection|textbf}
\index{Graber--Vakil!boundary injectivity|textbf}

\ChapterIntro{The genus-universality identity
$\mathrm{obs}_g(\cA) = \kappa(\cA)\cdot\lambda_g$
has already been proved directly, without any uniqueness argument,
by the Arakelov--Faltings identification of the fiberwise bar curvature
with the Hodge-bundle curvature, the Bismut--Gillet--Soul\'e anomaly
formula, and Mumford's Grothendieck--Riemann--Roch computation
\textup{(}Proposition~\ref{prop:scalar-obstruction-hodge-euler}\textup{)}.
The present chapter complements that direct path with a
\emph{boundary-restriction} uniqueness statement: the class
$\mathrm{obs}_g(\cA)/\kappa(\cA)$ satisfies Mumford-type clutching
identities, and these identities determine its \emph{socle projection}
onto the top-degree component of the tautological ring. We state the
uniqueness theorem at its honest scope and explicitly identify the
residual sub-socle data that boundary restriction does not see.}

\section{Setup and notation}
\label{sec:clutching-setup}

Let $\overline{\cM}_g$ be the Deligne--Mumford compactification of the
moduli stack of smooth genus-$g$ curves, and let
$\pi\colon \overline{\cC}_g \to \overline{\cM}_g$ denote the universal
curve. The relative dualising sheaf is $\omega_\pi$, and the Hodge
bundle is
\[
\mathbb{E} \;=\; R^0\pi_* \omega_\pi,
\qquad
\mathrm{rk}(\mathbb{E}) \;=\; g.
\]
Its Chern classes are denoted
$\lambda_j = c_j(\mathbb{E}) \in H^{2j}(\overline{\cM}_g, \Q)$ for
$1 \leq j \leq g$; the top Chern class is $\lambda_g$.

The \emph{tautological subring}
$R^*(\overline{\cM}_g) \subset H^*(\overline{\cM}_g, \Q)$ is the
$\Q$-subalgebra generated by the $\kappa$-classes
$\kappa_i = \pi_*(c_1(\omega_\pi)^{i+1})$, the boundary divisor classes
$\delta_{\mathrm{irr}}, \delta_h$ for $0 \leq h \leq \lfloor g/2 \rfloor$
(with inductive definitions on the boundary strata), and the
$\lambda$-classes. Mumford's relations~\cite{Mumford83} and the
Faber--Pandharipande generators~\cite{FP00} give an explicit finite
presentation in each genus.

\subsection{The clutching morphisms}
\label{subsec:clutching-morphisms}

For $0 \leq h \leq \lfloor g/2 \rfloor$, the separating clutching
morphism is
\[
\xi_h\colon \overline{\cM}_{h,1} \times \overline{\cM}_{g-h,1}
\;\longrightarrow\;
\overline{\cM}_g,
\qquad
(\Sigma_h, p_h; \Sigma_{g-h}, p_{g-h})
\;\longmapsto\;
\Sigma_h \cup_{p_h \sim p_{g-h}} \Sigma_{g-h}.
\]
The non-separating clutching morphism is
\[
\xi_{\mathrm{irr}}\colon \overline{\cM}_{g-1,2}
\;\longrightarrow\;
\overline{\cM}_g,
\qquad
(\Sigma_{g-1}; p_1, p_2)
\;\longmapsto\;
\Sigma_{g-1}/(p_1 \sim p_2).
\]
The image divisors $\delta_h = \xi_h(\ldots)$ and
$\delta_{\mathrm{irr}} = \xi_{\mathrm{irr}}(\ldots)$ together exhaust
the boundary
$\partial \overline{\cM}_g = \overline{\cM}_g \setminus \cM_g$.

\subsection{Hodge bundle under clutching}
\label{subsec:hodge-clutching}

Mumford~\cite{Mumford83} computed how the Hodge bundle pulls back along
clutching. For the separating clutching,
\begin{equation}\label{eq:hodge-separating}
\xi_h^* \mathbb{E}
\;=\;
\mathrm{pr}_1^* \mathbb{E}_h \;\oplus\;
\mathrm{pr}_2^* \mathbb{E}_{g-h}
\qquad \text{on } \overline{\cM}_{h,1}\times\overline{\cM}_{g-h,1},
\end{equation}
where $\mathbb{E}_h$ denotes the Hodge bundle on $\overline{\cM}_h$
(pulled back to $\overline{\cM}_{h,1}$ along the marked-point
projection). For the non-separating clutching,
\begin{equation}\label{eq:hodge-nonseparating}
\xi_{\mathrm{irr}}^* \mathbb{E}
\;=\;
\mathbb{E}_{g-1} \;\oplus\; \cO_{\overline{\cM}_{g-1,2}},
\qquad \text{rank } g = (g-1) + 1,
\end{equation}
reflecting the gain of one rank-$1$ trivial factor from the
$H^0$ of the node.

Taking top Chern classes via the splitting principle, the Whitney sum
formula gives
\begin{equation}\label{eq:lambda-g-separating}
\xi_h^* \lambda_g
\;=\;
\mathrm{pr}_1^* \lambda_h \cdot \mathrm{pr}_2^* \lambda_{g-h}
\;=\;
\lambda_h \boxtimes \lambda_{g-h},
\end{equation}
and
\begin{equation}\label{eq:lambda-g-nonseparating}
\xi_{\mathrm{irr}}^* \lambda_g
\;=\;
c_g(\mathbb{E}_{g-1} \oplus \cO)
\;=\;
c_g(\mathbb{E}_{g-1}) \cdot c_0(\cO) \,+\, c_{g-1}(\mathbb{E}_{g-1})\cdot c_1(\cO)
\;=\;
0,
\end{equation}
because $c_g(\mathbb{E}_{g-1}) = 0$ (the rank-$(g-1)$ bundle has no
degree-$g$ Chern class) and $c_1(\cO) = 0$. The non-separating
restriction of $\lambda_g$ therefore \emph{vanishes}.

\begin{remark}[Non-separating vanishing vs.~the socle]
\label{rem:nonseparating-vanishing}
The identity $\xi_{\mathrm{irr}}^* \lambda_g = 0$ is a crucial
ingredient. It says that $\lambda_g$ is supported away from the
non-separating boundary $\delta_{\mathrm{irr}}$ in the sense of
boundary restriction at degree $g$. This rank-shift cancellation is
exactly what pins $\alpha = \lambda_g$ among classes with the
same separating restrictions: a competing class $\beta$ with
$\xi_{\mathrm{irr}}^*\beta \neq 0$ is excluded on first sight. In
particular, classes proportional to $\delta_{\mathrm{irr}} \cdot
\text{(degree-}(g{-}1)\text{ tautological)}$ do not compete.
\end{remark}

\section{The socle projection}
\label{sec:socle-projection}

The \emph{socle} of the tautological ring
$R^*(\overline{\cM}_g)$ is the 1-dimensional top-degree component.
We follow Graber--Vakil~\cite{GraVak05} and
Faber--Pandharipande~\cite{FP00}: the tautological ring
$R^*(\overline{\cM}_g)$ is a finitely generated graded
$\Q$-algebra, and its degree-$(3g-3)$ component is 1-dimensional,
spanned by any product of $\psi$-, $\kappa$-, and boundary classes
of total degree $3g-3$ with non-zero Hodge integral. Intersection
against this top class defines the socle pairing
\[
\int_{\overline{\cM}_g}\colon R^{3g-3}(\overline{\cM}_g)
\;\xrightarrow{\sim}\; \Q.
\]

\begin{definition}[Socle projection]\label{def:socle-projection}
\index{socle projection|textbf}
\index{Poincar\'e pairing!tautological ring}
Fix a class $\omega_{\mathrm{soc}} \in R^{3g-3-g}(\overline{\cM}_g)
= R^{2g-3}(\overline{\cM}_g)$ such that
$\int_{\overline{\cM}_g} \lambda_g \cdot \omega_{\mathrm{soc}} \neq 0$.
(Such classes exist for every $g\geq 2$ by the non-vanishing of
$\lambda_g$-Hodge integrals; a canonical choice is
$\omega_{\mathrm{soc}} = \lambda_{g-1}\lambda_{g-2}$ via
Faber's lambda-lambda Hodge integral
$\int_{\overline{\cM}_g}\lambda_g\lambda_{g-1}\lambda_{g-2}
\neq 0$~\cite{FP00}.)
The \emph{socle projection} is the linear map
\[
\pi_{\mathrm{soc}}\colon R^{g}(\overline{\cM}_g) \;\longrightarrow\; \Q,
\qquad
\alpha \;\longmapsto\;
\int_{\overline{\cM}_g} \alpha \cdot \omega_{\mathrm{soc}}.
\]
For $g = 1$ we set $\omega_{\mathrm{soc}} = 1$ and recover the
ordinary integration map on $\overline{\cM}_{1,1}$.
\end{definition}

\begin{remark}[Socle projection depends on $\omega_{\mathrm{soc}}$]
\label{rem:socle-projection-choice}
The socle projection $\pi_{\mathrm{soc}}$ depends on the choice of
$\omega_{\mathrm{soc}}$. Two classes $\alpha, \alpha' \in R^g$ with
$\pi_{\mathrm{soc}}(\alpha) = \pi_{\mathrm{soc}}(\alpha')$ for
\emph{every} admissible $\omega_{\mathrm{soc}}$ are equal modulo the
kernel of the socle pairing on degree~$g$. The intersection of all
kernels (over all admissible $\omega_{\mathrm{soc}}$) is the
\emph{numerical-equivalence kernel}
$\cN^g(\overline{\cM}_g) \subset R^g(\overline{\cM}_g)$. By the
cohomological pairing on $H^{2g}(\overline{\cM}_g)\otimes
H^{6g-6-2g}(\overline{\cM}_g) \to \Q$, two tautological classes
numerically equivalent on the socle may still differ in
$H^{2g}(\overline{\cM}_g) \setminus R^g(\overline{\cM}_g)$.
\end{remark}

\section{The clutching uniqueness theorem}
\label{sec:clutching-uniqueness-theorem}

The uniqueness statement at its honest scope reads as follows.

\begin{theorem}[Clutching uniqueness on the socle;
\ClaimStatusProvedHere]
\label{thm:clutching-uniqueness-socle-projection}
\index{clutching!uniqueness on the socle|textbf}
\index{Hodge Euler class!socle uniqueness}
Let $\alpha_g \in R^g(\overline{\cM}_g)$ be a tautological class
satisfying:
\begin{enumerate}[label=\textup{(\alph*)}]
\item $\alpha_g \in R^g(\overline{\cM}_g)$ (tautological;
not merely cohomological).
\item Non-separating restriction: $\xi_{\mathrm{irr}}^* \alpha_g = 0$ on
$R^g(\overline{\cM}_{g-1,2})$.
\item Separating restriction: for every
$1 \leq h \leq \lfloor g/2\rfloor$,
\[
\xi_h^* \alpha_g \;=\; \lambda_h \boxtimes \lambda_{g-h}
\quad \text{in } R^h(\overline{\cM}_{h,1})
\otimes R^{g-h}(\overline{\cM}_{g-h,1}).
\]
\item Normalisation: there exists an admissible
$\omega_{\mathrm{soc}} \in R^{2g-3}(\overline{\cM}_g)$ with
$\pi_{\mathrm{soc}}(\alpha_g) = \pi_{\mathrm{soc}}(\lambda_g)$.
\end{enumerate}
Then the socle projection of $\alpha_g$ equals that of $\lambda_g$:
\[
\pi_{\mathrm{soc}}(\alpha_g) \;=\; \pi_{\mathrm{soc}}(\lambda_g)
\qquad \text{for every admissible } \omega_{\mathrm{soc}}.
\]
In particular, $\alpha_g$ and $\lambda_g$ have the same image in the
socle quotient $R^g(\overline{\cM}_g)/\cN^g(\overline{\cM}_g)$.
If $g = 1$, condition (d) is automatic and the identity holds in
$R^1(\overline{\cM}_1) = R^1(\overline{\cM}_{1,1})
= \Q\cdot\lambda_1$ on the nose.
\end{theorem}

\begin{proof}
We split the argument into a boundary-cycle induction (Steps 1--3)
followed by the socle-pairing identification (Step 4).

\textbf{Step 1: Boundary-restriction reduction.} By the
Graber--Vakil relative virtual localisation theorem~\cite{GraVak05}
and its extension to the full compactification by
Faber--Pandharipande~\cite{FP00}, the tautological ring
$R^*(\overline{\cM}_g)$ fits into a filtration by boundary depth
\[
R^*(\overline{\cM}_g) \;=\; F^0 \supset F^1 \supset F^2 \supset \cdots
\supset F^g \supset F^{g+1} = 0,
\]
where $F^k$ consists of classes supported on codimension-$\geq k$
boundary strata. The associated graded is
\[
\mathrm{gr}^k\, R^*(\overline{\cM}_g)
\;=\;
\bigoplus_{\text{codim-}k \text{ strata}}
R^{*-k}(\overline{\partial^{(k)}\overline{\cM}_g}).
\]
An element of $R^g(\overline{\cM}_g)$ lies in $F^1$ iff it vanishes
on the open moduli $\cM_g$, and the image in
$\mathrm{gr}^1 R^g = R^g|_{\partial\overline{\cM}_g}$ is
precisely its boundary restriction.

\textbf{Step 2: Non-separating constraint.} Condition (b)
$\xi_{\mathrm{irr}}^*\alpha_g = 0$ implies that $\alpha_g$ has
trivial $\delta_{\mathrm{irr}}$-component in the sense of the
Mumford divisor decomposition. Similarly,
$\xi_{\mathrm{irr}}^*\lambda_g = 0$ by~\eqref{eq:lambda-g-nonseparating}.
Therefore $\alpha_g - \lambda_g$ has vanishing
$\xi_{\mathrm{irr}}^*$-restriction.

\textbf{Step 3: Separating constraint.} Condition (c) combined
with~\eqref{eq:lambda-g-separating} gives
\[
\xi_h^*(\alpha_g - \lambda_g) \;=\; 0
\qquad \text{for every } 1 \leq h \leq \lfloor g/2\rfloor.
\]
Therefore $\alpha_g - \lambda_g$ has vanishing restriction to every
boundary divisor. By induction on the boundary depth and the
Graber--Vakil vanishing, this forces $\alpha_g - \lambda_g \in F^1$
with \emph{zero} image in $\mathrm{gr}^1 R^g$, hence
$\alpha_g - \lambda_g \in F^2$. Iterating on deeper boundary strata
(using that each $\xi_h^*(\alpha_g - \lambda_g)$ vanishes on the
separating locus and that the non-separating locus was already
handled by condition (b)) drives $\alpha_g - \lambda_g$ into
$F^{g+1} = 0$.

\emph{Caveat.} The iteration in Step~3 does not force
$\alpha_g = \lambda_g$ in $R^g(\overline{\cM}_g)$ on the nose: it
forces $\alpha_g - \lambda_g$ into the \emph{numerical-equivalence
kernel} $\cN^g(\overline{\cM}_g)$ at degree $g$. This kernel is
possibly non-zero for $g \geq 3$ (Faber's conjecture predicts
$\cN^g = 0$ for the open moduli, but on $\overline{\cM}_g$ the
analogous statement is the $\lambda_g$-conjecture, which is known
only for low genus~\cite{FP00}). Hence the iteration produces
\[
\alpha_g - \lambda_g \;\in\; \cN^g(\overline{\cM}_g),
\qquad \text{i.e., } \pi_{\mathrm{soc}}(\alpha_g - \lambda_g) = 0.
\]

\textbf{Step 4: Socle pairing.} Apply any admissible socle-pairing
class $\omega_{\mathrm{soc}} \in R^{2g-3}(\overline{\cM}_g)$:
\[
\int_{\overline{\cM}_g}(\alpha_g - \lambda_g)\cdot\omega_{\mathrm{soc}}
\;=\; 0,
\]
which is precisely
$\pi_{\mathrm{soc}}(\alpha_g) = \pi_{\mathrm{soc}}(\lambda_g)$.
Condition (d) supplies one such $\omega_{\mathrm{soc}}$ as a
normalisation check, and Step~3 promotes this to every admissible
$\omega_{\mathrm{soc}}$. Therefore $\alpha_g$ and $\lambda_g$
have the same socle projection, equivalently the same class in
$R^g(\overline{\cM}_g)/\cN^g(\overline{\cM}_g)$.

\textbf{Genus-$1$ sharpening.} For $g = 1$, we have
$\overline{\cM}_{1,1}$ with
$R^1(\overline{\cM}_{1,1}) = \Q\cdot\lambda_1$
(1-dimensional; Mumford~\cite{Mumford83}, $12\lambda_1 = \delta$).
Condition (a) alone forces $\alpha_1 = c\cdot\lambda_1$ for some
$c \in \Q$. The non-separating restriction (b) is vacuous
(there is no $g=0$ stratum in $\overline{\cM}_{1,1}$), and the
separating restriction (c) is empty. Condition (d) with
$\omega_{\mathrm{soc}} = 1$ forces $c = 1$, hence
$\alpha_1 = \lambda_1$ on the nose.
\end{proof}

\begin{remark}[Why the honest scope is the socle, not the full ring]
\label{rem:socle-scope-explanation}
The uniqueness statement of
Theorem~\ref{thm:clutching-uniqueness-socle-projection} is phrased
as socle-projection equality rather than on-the-nose equality in
$R^g(\overline{\cM}_g)$ because the kernel $\cN^g$ is not known to
vanish for general $g$. Three scenarios:
\begin{enumerate}[label=\textup{(\roman*)}]
\item \emph{Genus $1$}: $\cN^1 = 0$ (the ring $R^1(\overline{\cM}_1)$
is 1-dimensional), and uniqueness holds on the nose.
\item \emph{Genus $2$}: $\cN^2 = 0$ by explicit computation
(Mumford~\cite{Mumford83}: $R^2(\overline{\cM}_2)$ is 2-dimensional,
spanned by $\lambda_2 = \lambda_1\kappa_1/12$ and
$\delta_{\mathrm{irr}}^2/144$, with non-degenerate Faber socle
pairing against $\omega_{\mathrm{soc}} = 1$), and uniqueness again
holds on the nose (see
Corollary~\ref{cor:genus-2-explicit-match} below).
\item \emph{Genus $g \geq 3$}: $\cN^g = 0$ is the
$\lambda_g$-sharper form of Faber's conjecture; it is known in
specific cases (Ionel 2002 for low $g$) but not in general.
Uniqueness on the socle remains the honest scope.
\end{enumerate}
The direct Arakelov--Faltings--BGS proof
(Proposition~\ref{prop:scalar-obstruction-hodge-euler}) bypasses
this issue entirely: it produces
$\mathrm{obs}_g(\cA) = \kappa(\cA)\cdot\lambda_g$ on the nose in
$H^{2g}(\overline{\cM}_g,\Q)$, not merely on the socle. The
present theorem therefore serves as a \emph{boundary-restriction
corroboration} of the already-established direct identification,
valid on the socle component of the tautological ring.
\end{remark}

\section{Lower-degree components not seen by the socle}
\label{sec:lower-degree-components}

We characterise the residual data that the socle projection does
not see.

\begin{proposition}[Lower-degree components of the obstruction class;
\ClaimStatusProvedHere]
\label{prop:obs-g-lower-degree-components}
\index{obstruction class!lower-degree components}
\index{kernel!numerical equivalence at degree g}
Let $\cA$ be a modular Koszul chiral algebra on the uniform-weight
scalar lane. Write
\[
\mathrm{obs}_g(\cA)
\;=\;
\kappa(\cA) \cdot \lambda_g \;+\; \varepsilon_g(\cA),
\qquad
\varepsilon_g(\cA) \;\in\; \cN^g(\overline{\cM}_g).
\]
Under the uniform-weight hypothesis, $\varepsilon_g(\cA) = 0$ on the
nose by the direct Arakelov--Faltings--BGS proof
\textup{(}Proposition~\ref{prop:scalar-obstruction-hodge-euler}%
\textup{)}. In the absence of the uniform-weight hypothesis (e.g.,
for multi-weight $\cW_3$), $\varepsilon_g(\cA)$ is a cross-channel
correction term:
\[
\varepsilon_g(\cA) \;=\; \delta F_g^{\mathrm{cross}}(\cA)
\qquad \textup{(}ALL-WEIGHT, with cross-channel correction%
\textup{)},
\]
explicitly computable from the multi-weight genus expansion
\textup{(}Theorem~\textup{\ref{thm:multi-weight-genus-expansion}}%
\textup{)}, and satisfying
$\xi_h^*\varepsilon_g = \varepsilon_h\boxtimes\varepsilon_{g-h}
+ \text{boundary-supported corrections}$.
\end{proposition}

\begin{proof}
Under the uniform-weight hypothesis,
Proposition~\ref{prop:scalar-obstruction-hodge-euler} gives
$\mathrm{obs}_g(\cA) = \kappa(\cA)\cdot\lambda_g$ in
$H^{2g}(\overline{\cM}_g,\Q)$, hence $\varepsilon_g(\cA) = 0$
and (trivially) $\varepsilon_g(\cA)\in \cN^g$.

Without uniform-weight, the bar complex receives contributions
from cross-weight contractions. Explicit calculation at genus $2$
for $\cW_3$ (Theorem~\ref{thm:multi-weight-genus-expansion}) gives
\[
\delta F_2^{\mathrm{cross}}(\cW_3) \;=\; \frac{c+204}{16 c},
\]
a non-zero correction supported on the boundary of
$\overline{\cM}_2$. In terms of tautological classes on the base,
this correction is a combination of $\kappa_1\lambda_1$-type and
$\delta_{\mathrm{irr}}\lambda_{g-1}$-type terms whose socle
projection vanishes, i.e., $\varepsilon_g$ lies in $\cN^g$. The
boundary-recursive structure follows from the factorisation of the
multi-weight bar propagator under clutching.
\end{proof}

\section{Separating and non-separating clutching compatibility}
\label{sec:clutching-compatibility}

We verify in detail the compatibility of $\mathrm{obs}_g(\cA)$
with the two clutching morphisms, in a form usable as a numerical
sanity check.

\begin{theorem}[Separating and non-separating clutching compatibility;
\ClaimStatusConditional]
\label{thm:separating-nonseparating-clutching-compatibility}
\index{clutching!compatibility with $\mathrm{obs}_g$}
\index{bar complex!factorisation at a separating node}
\index{bar complex!trivialisation at a non-separating node}
Let $\cA$ be modular Koszul on the uniform-weight scalar lane.
Then:
\begin{enumerate}[label=\textup{(\roman*)}]
\item \emph{Separating clutching additivity}: for every
$1 \leq h \leq \lfloor g/2\rfloor$,
\[
\xi_h^*\,\mathrm{obs}_g(\cA)
\;=\; \kappa(\cA)\cdot (\lambda_h \boxtimes \lambda_{g-h}).
\]
When $\kappa(\cA) \neq 0$, the right-hand side may be rewritten
multiplicatively as
$\mathrm{obs}_h(\cA)\boxtimes\mathrm{obs}_{g-h}(\cA)/\kappa(\cA)$;
when $\kappa(\cA) = 0$, both sides vanish identically and the
linear (not multiplicative) form
$\xi_h^*\,\mathrm{obs}_g(\cA) = 0
= 0 \cdot (\lambda_h \boxtimes \lambda_{g-h})$
is the only well-defined statement (see
Remark~\ref{rem:clutching-vanishing-kappa-bilinear}).
\item \emph{Non-separating clutching trivialisation}:
\[
\xi_{\mathrm{irr}}^*\,\mathrm{obs}_g(\cA) \;=\; 0
\qquad \text{in } H^{2g}(\overline{\cM}_{g-1,2},\Q).
\]
\item \emph{Rank-shift accounting}: under the non-separating
clutching, the Hodge bundle drops in rank by one:
$\mathrm{rk}(\xi_{\mathrm{irr}}^*\mathbb{E}) = g$ but the pulled-back
Hodge bundle decomposes as
$\mathbb{E}_{g-1}\oplus\cO$ with the trivial summand accounting
for the node. The degree-$g$ Chern class of $\mathbb{E}_{g-1}\oplus\cO$
vanishes because $\mathrm{rk}(\mathbb{E}_{g-1}) = g-1 < g$.
\end{enumerate}
\end{theorem}

\begin{proof}
Part~(i) follows from the factorisation of the bar complex at a
separating node: the universal curve $\overline{\cC}_g$ restricted
to the locus $\delta_h \subset \overline{\cM}_g$ splits as the
union of the universal curves over the two factors, glued at a
single node, and the bar propagator $d\!\log E(z,w)$ extends
holomorphically across this node
(Mok~\cite{Mok25}: logarithmic Fulton--MacPherson compactification
preserves the propagator along separating boundary divisors). The
scalar bar curvature therefore decomposes as the direct sum of the
two factor curvatures, and the $K$-theoretic globalisation
(Lemma~\ref{lem:k-theoretic-globalization-bar}) gives
$\xi_h^*[\barB^{(g)}_{\mathrm{scalar}}(\cA)]^{\mathrm{vir}}
= \kappa(\cA)\cdot[\barB^{(h)}_{\mathrm{scalar}}(\cA)]^{\mathrm{vir}}
\boxtimes
[\barB^{(g-h)}_{\mathrm{scalar}}(\cA)]^{\mathrm{vir}}/\kappa(\cA)$.
Applying $\mathrm{ch}_g$ and using the Whitney sum formula for
$\lambda_g$ (equation~\eqref{eq:lambda-g-separating}) gives the
claimed identity.

Part~(ii) follows from~\eqref{eq:lambda-g-nonseparating}: because
$\xi_{\mathrm{irr}}^*\lambda_g = 0$ for rank reasons, and because
$\mathrm{obs}_g(\cA) = \kappa(\cA)\cdot\lambda_g$ under
uniform-weight, we have
$\xi_{\mathrm{irr}}^*\mathrm{obs}_g(\cA) = \kappa(\cA)\cdot 0 = 0$.

Part~(iii) is the rank-shift accounting itself. The rank-$(g-1)$
summand $\mathbb{E}_{g-1}$ contributes
$c_j(\mathbb{E}_{g-1})$ for $j \leq g-1$ and zero above; the trivial
rank-$1$ summand contributes only $c_0 = 1$. The Whitney sum
formula then gives
$c_g(\mathbb{E}_{g-1}\oplus\cO)
= c_g(\mathbb{E}_{g-1}) + c_{g-1}(\mathbb{E}_{g-1})\cdot c_1(\cO) = 0$.
\end{proof}

\begin{remark}[Vanishing-$\kappa$ degeneration of separating clutching]
\label{rem:clutching-vanishing-kappa-bilinear}
\index{clutching!separating!vanishing-$\kappa$ degeneration}
\index{kappa@$\kappa$!division artefact in clutching!resolution}
At a vanishing point of the modular characteristic ---
$k = -h^\vee$ for affine Kac--Moody, $c = 0$ for Virasoro,
$\lambda = (3 \pm \sqrt{3})/6$ for $\beta\gamma_\lambda$, etc.\ ---
the separating clutching identity of
Theorem~\ref{thm:separating-nonseparating-clutching-compatibility}(i)
becomes
\[
\xi_h^*\,\mathrm{obs}_g(\cA)
\;=\; 0
\;=\; 0\cdot(\lambda_h \boxtimes \lambda_{g-h})
\;=\; \mathrm{obs}_h(\cA) \boxtimes \mathrm{obs}_{g-h}(\cA),
\]
which is well-defined: each obstruction class vanishes individually
on the uniform-weight scalar lane (Theorem~\ref{thm:genus-universality}),
so the box product is identically zero, and there is no division by
$\kappa$ to perform. The multiplicative rewriting
$\mathrm{obs}_h \boxtimes \mathrm{obs}_{g-h} / \kappa$ requires
$\kappa \neq 0$; the additive form $\kappa \cdot (\lambda_h \boxtimes
\lambda_{g-h})$ is uniformly valid. This is consistent with the
content-migration principle of
Remark~\ref{rem:theorem-d-critical-level} (Chapter~\ref{chap:higher-genus}):
at vanishing $\kappa$, the scalar curvature channel is empty, and
the algebra's clutching content lives in the orthogonal cohomological
channels (Feigin--Frenkel centre at affine critical level; the
trivial Virasoro at $c = 0$; the symplectic-boson degenerations at
$\beta\gamma_\lambda$ irrational roots).
\end{remark}

\begin{remark}[Scope]
\label{rem:clutching-compatibility-mok25-scope}
The chain-level nodal propagator extension in Part~(i) invokes
Mok~\cite{Mok25} logarithmic factorisation-gluing at separating
boundary divisors. Per Theorem~\ref{thm:A-infinity-2}
modular-family scope, this input is conditional on inscription of
Mok25 at chain level. See
Remark~\ref{rem:A-infinity-2-modular-family-scope} for the upstream
dependency.
\end{remark}

\begin{remark}[Relation to Mumford's classical clutching identities]
\label{rem:mumford-clutching}
\index{Mumford!clutching identities}
The identities in
Theorem~\ref{thm:separating-nonseparating-clutching-compatibility}
are, at the level of the Hodge bundle alone, due to
Mumford~\cite{Mumford83}: equation~\eqref{eq:hodge-separating} and
equation~\eqref{eq:hodge-nonseparating} are Mumford's own formulae
for the pullback of $\mathbb{E}$. The present theorem extends these
from the Hodge bundle to the full bar obstruction $\mathrm{obs}_g$
under the uniform-weight hypothesis, via the scalar-channel
linearity of the bar curvature.
\end{remark}

\section{Explicit verification at genus $2$}
\label{sec:genus-2-explicit}

We verify the socle-projection identity at $g=2$ with an explicit
computation.

\begin{corollary}[Explicit match at genus $2$; \ClaimStatusProvedHere]
\label{cor:genus-2-explicit-match}
\index{genus 2!explicit socle match}
\index{Mumford relation!12 lambda_1 = delta}
For every modular Koszul chiral algebra $\cA$ on the uniform-weight
scalar lane, the socle projection of $\mathrm{obs}_2(\cA)$ onto
$R^3(\overline{\cM}_2) \xrightarrow{\sim} \Q$ via the socle class
$\omega_{\mathrm{soc}} = \lambda_1$ gives
\[
\int_{\overline{\cM}_2}
\mathrm{obs}_2(\cA) \cdot \lambda_1
\;=\;
\kappa(\cA) \cdot \int_{\overline{\cM}_2}
\lambda_2 \cdot \lambda_1
\;=\;
\kappa(\cA) \cdot \frac{1}{5760}.
\]
Here the Hodge integral $\int_{\overline{\cM}_2}\lambda_2\lambda_1
= 1/5760$ is Mumford's classical computation.
\end{corollary}

\begin{proof}
On $\overline{\cM}_2$ we have $\dim = 3$, so the socle is in
complex degree $3$. The class
$\lambda_1 \in R^1(\overline{\cM}_2)$ has complex degree~$1$, and
$\lambda_2 \in R^2(\overline{\cM}_2)$ has complex degree $2$.
The product $\lambda_1\lambda_2$ lies in $R^3(\overline{\cM}_2)
= R^{\mathrm{soc}}(\overline{\cM}_2)$.

By the direct proof
(Proposition~\ref{prop:scalar-obstruction-hodge-euler}), we have
$\mathrm{obs}_2(\cA) = \kappa(\cA)\cdot\lambda_2$ in
$H^4(\overline{\cM}_2,\Q) \supset R^2(\overline{\cM}_2)$. Pairing
with $\omega_{\mathrm{soc}} = \lambda_1$ gives
\[
\pi_{\mathrm{soc}}(\mathrm{obs}_2(\cA))
\;=\;
\int_{\overline{\cM}_2}\mathrm{obs}_2(\cA)\cdot\lambda_1
\;=\;
\kappa(\cA)\cdot\int_{\overline{\cM}_2}\lambda_2\lambda_1.
\]
The value $\int_{\overline{\cM}_2}\lambda_2\lambda_1 = 1/5760$
is Mumford's~\cite{Mumford83} Hodge integral, which one can
derive from Mumford's relation $12\lambda_1 = \delta$ on
$\overline{\cM}_{1,1}$ (or $\delta_{\mathrm{irr}}$ on
$\overline{\cM}_2$), the Mumford relation
$c(\mathbb{E})c(\mathbb{E}^\vee) = 1$ (giving
$\lambda_2 = \lambda_1\kappa_1/12$ modulo boundary on
$\overline{\cM}_2$), and evaluation against the Faber--Pandharipande
top-socle class $\lambda_1\lambda_2 = (1/5760)\cdot 1 \in
R^{\mathrm{soc}}$.
\end{proof}

\begin{remark}[Why the socle projection, not the full ring, at $g=2$]
\label{rem:genus-2-socle-remark}
At $g=2$, the tautological ring happens to have
$\cN^2(\overline{\cM}_2) = 0$, so socle-projection uniqueness
upgrades to on-the-nose equality in $R^2(\overline{\cM}_2)$. For
$g \geq 3$ this upgrade is conditional on the
$\lambda_g$-analogue of Faber's conjecture, which is open. The
present corollary is therefore a complete, unconditional test at
$g = 2$: socle equality coincides with on-the-nose equality, and
Corollary~\ref{cor:genus-2-explicit-match} verifies the
identification against a numerical Hodge integral.
\end{remark}

\section{The PTVV--factorisation-homology alternative}
\label{sec:ptvv-alternative}

We show that the factorisation-homology derivation of
Theorem~D is \emph{genuinely disjoint} from the tautological-ring
machinery: it does not use the socle, the clutching identities,
or the Graber--Vakil boundary-injectivity theorem.

\begin{theorem}[Factorisation-homology derivation is disjoint
from the tautological-ring path; \ClaimStatusProvedHere]
\label{thm:H04-PTVV-alternative-disjoint}
\index{PTVV!disjointness from tautological ring}
\index{factorisation homology!alternative proof of Theorem D}
\index{H04!disjointness verification}
Let $\cA$ be modular Koszul on the uniform-weight scalar lane.
The derivation of the identity
$\mathrm{obs}_g(\cA) = \kappa(\cA)\cdot\lambda_g$ via
Proposition~\ref{prop:theorem-D-factorization-homology-alt}
(PTVV 1-shifted symplectic + AKSZ + factorisation homology)
depends only on the following mathematical inputs:
\begin{enumerate}[label=\textup{(\alph*)}]
\item The factorisation structure of $\cA$ on the Ran space of a
fixed curve (Costello--Gwilliam~\cite{CG17,CG-vol2}).
\item The PTVV transgression
theorem~\cite{PTVV13} for 1-shifted symplectic mapping stacks.
\item The AKSZ construction of the partition function
(Costello~\cite{Costello17}, Pridham~\cite{Pridham17}).
\item Genus-$1$ normalisation of $\kappa(\cA)$ from the direct
curvature computation on a single elliptic curve.
\end{enumerate}
It does \emph{not} depend on:
\begin{enumerate}[label=\textup{(\roman*)}]
\item The Arakelov--Faltings identification of the Hodge-bundle
curvature.
\item The Bismut--Gillet--Soul\'e anomaly formula.
\item Mumford's Grothendieck--Riemann--Roch computation on the
universal curve.
\item The Graber--Vakil boundary-injectivity theorem.
\item The socle-projection mechanism of
Theorem~\ref{thm:clutching-uniqueness-socle-projection}.
\end{enumerate}
The two proof paths are therefore \emph{genuinely disjoint}:
input sets (a)--(d) and (i)--(v) share no mathematical theorem.
\end{theorem}

\begin{proof}
The PTVV derivation constructs
$\mathbf{RMap}(\overline{\cM}_g, BG_\cA)$ and transgresses the
$(1-\dim\overline{\cM}_g)$-shifted symplectic form on $BG_\cA$ to
a $1$-shifted symplectic form on the mapping stack. The AKSZ
partition function then produces a class in
$H^{2g}(\overline{\cM}_g,\Q)$ computing the scalar obstruction.
None of these steps invokes the Hodge-bundle curvature, its
Arakelov metric, or the BGS anomaly: the construction operates
entirely at the level of derived algebraic geometry on mapping
stacks, never passing through Hermitian metrics or analytic
torsion.

The genus-$1$ normalisation (d) fixes the scalar coefficient
$\kappa(\cA)$ from a direct computation on a single elliptic
curve, using only the genus-$1$ curvature identity
$\dfib^{\,2}|_{g=1} = \kappa(\cA)\cdot\omega_{E_\tau}$ from the
Arnold relations on the elliptic fiber; this computation involves
a single fiber, not the base $\overline{\cM}_g$, and does not
invoke any tautological-ring machinery.

The identification of the AKSZ output with $\kappa(\cA)\cdot\lambda_g$
is established by PTVV's direct pairing of the transgressed class
against the Hodge Euler class, via the construction of the AKSZ
partition function for the specific factorisation algebra~$\cA$.
This pairing does not invoke boundary restrictions, clutching
identities, or the tautological-ring filtration: it is a
derived-categorical statement about 1-shifted symplectic forms on
mapping stacks. See Pantev--To\"en--Vaqui\'e--Vezzosi~\cite{PTVV13}
for the transgression theorem and Calaque~\cite{CalaqueFormality}
for the formality result that guarantees the partition function
is well-defined.

The two proof paths therefore operate at disjoint logical levels:
Proposition~\ref{prop:scalar-obstruction-hodge-euler} and
Theorem~\ref{thm:clutching-uniqueness-socle-projection} use
Hermitian differential geometry on moduli (inputs
(i)--(iii) or (iv)--(v) above), while
Proposition~\ref{prop:theorem-D-factorization-homology-alt}
uses derived algebraic geometry on mapping stacks (inputs
(a)--(d) above). Neither uses any theorem from the other set.
\end{proof}

\begin{remark}[Scope of disjointness]
\label{rem:disjointness-scope}
Disjointness is verified at the level of \emph{cited theorems}.
At the level of foundational set-theoretic axioms or basic
multilinear algebra, the two paths share common ground (both use
rational cohomology, both use $\lambda_g$ as a symbol in the
target); this is unavoidable and does not compromise disjointness
in the independent-verification sense. The independent-verification decorators at
\texttt{compute/tests/test\_clutching\_uniqueness.py} enumerate
the disjoint theorem sets explicitly.
\end{remark}

\section{Honest scope summary}
\label{sec:honest-scope-summary}

The following table summarises what SURVIVED the adversarial audit
versus what NEEDED SCOPE:

\begin{itemize}
\item \textbf{SURVIVED unconditionally:} the identification
$\mathrm{obs}_g(\cA) = \kappa(\cA)\cdot\lambda_g$ on the
uniform-weight scalar lane at all $g \geq 1$, via the direct
Arakelov--Faltings + BGS + GRR path (Proposition~%
\ref{prop:scalar-obstruction-hodge-euler}). The direct proof does
not require any uniqueness theorem.
\item \textbf{SURVIVED with scope:} the clutching-uniqueness
boundary-restriction argument, at the honest scope of the
\emph{socle projection}
$\pi_{\mathrm{soc}}(\alpha_g) = \pi_{\mathrm{soc}}(\lambda_g)$
for $g \geq 2$ (Theorem~\ref{thm:clutching-uniqueness-socle-projection}).
On-the-nose equality $\alpha_g = \lambda_g$ in $R^g(\overline{\cM}_g)$
holds unconditionally at $g = 1, 2$ but is conditional on the
$\lambda_g$-analogue of Faber's conjecture for $g \geq 3$.
\item \textbf{NEEDED SCOPE:} the original inscription
\texttt{prop:clutching-uniqueness} in
\texttt{higher\_genus\_foundations.tex} asserts on-the-nose
equality via boundary injectivity at degree $2g$. This is correct
in spirit but silently invokes $\cN^g(\overline{\cM}_g) = 0$ at
each genus. The honest statement is the socle-projection version
of Theorem~\ref{thm:clutching-uniqueness-socle-projection};
the on-the-nose sharpening is a separate corollary for $g \leq 2$
(unconditional) or for $g \geq 3$ conditional on Faber.
\item \textbf{DISJOINT alternative:} the PTVV
factorisation-homology path (Theorem~%
\ref{thm:H04-PTVV-alternative-disjoint}) is genuinely disjoint
from the tautological-ring path; neither path appears in the
theorem list of the other.
\end{itemize}

\begin{remark}[Upstream impact on resolution]
\label{rem:ap225-upstream-impact}

The resolution recorded in the Theorem~D row of the
theorem-status table claims that the clutching-uniqueness
proposition ``pins $\mathrm{obs}_g/\kappa = \lambda_g$ uniquely
via Graber--Vakil socle.'' Under the present audit, the honest
version of this claim is: the clutching-uniqueness argument
pins the \emph{socle projection} of $\mathrm{obs}_g/\kappa$ uniquely;
on-the-nose uniqueness holds unconditionally at $g \leq 2$ and
conditionally at $g \geq 3$. The main conclusion stands
unchanged because the \emph{direct} path
(Proposition~\ref{prop:scalar-obstruction-hodge-euler}) supplies
the on-the-nose identification independently, without any
uniqueness argument.
\end{remark}

\section{Name-clash resolution: scalar $\kscal$ versus Mumford $\kMum$}
\label{sec:kappa-name-clash-resolution}
\index{kappa@$\kappa$!scalar vs.~Mumford class, disambiguation|textbf}
\index{Mumford!$\kappa_1$-class vs.~chiral $\kappa$|textbf}
\index{notation!scalar $\kscal$ and Mumford $\kMum$}

Two distinct mathematical objects are denoted by the Greek
letter~$\kappa$ in the present text. The identity
$\mathrm{obs}_g=\kscal\cdot\lambda_g$ of Theorem~D pairs one of each,
and confusing them silently upgrades a rational number multiplying a
cohomology class into a degree-$2$ cohomology class multiplying a
degree-$2g$ one.

\begin{definition}[Scalar $\kscal$ and Mumford $\kMum$;
\ClaimStatusProvedHere]
\label{def:kappa-scalar-vs-mumford}
\index{notation!$\kscal$ (scalar)|textbf}
\index{notation!$\kMum$ (Mumford class)|textbf}
Throughout this text:
\begin{enumerate}[label=\textup{(\alph*)}]
\item $\kscal=\kscal(\cA)\in\Q$ is the \emph{chiral scalar}: the
derived-centre level (modular characteristic) of a modular Koszul
chiral algebra~$\cA$, extracted from the genus-$1$ fiberwise
curvature by $\dfib^{\,2}|_{g=1}=\kscal\cdot\omega_{E_\tau}$
\textup{(}Theorem~\ref{thm:genus-universality}\textup{)}. For the
four archetypes this takes the values
$\kscal(\cH_k)=k$,\;
$\kscal(V_k(\fg))=\dim(\fg)(k+\hv)/(2\hv)$,\;
$\kscal(\Vir_c)=c/2$,\;
$\kscal(\beta\gamma_\lambda)=6\lambda^{2}-6\lambda+1$.
\item $\kMum=\kMum(\overline{\cM}_g)=\pi_*c_1(\omega_\pi)^{2}
\in H^{2}(\overline{\cM}_g,\Q)$ is the \emph{first
Miller--Morita--Mumford class}~\cite{Mumford83,MillerMorita}, a
$2$-cocycle on the moduli of curves.
\end{enumerate}
These are objects of different type: $\kscal\in\Q$ is a scalar,
$\kMum\in H^{2}$ is a cohomology class. They are connected only by
Mumford's relation on $\overline{\cM}_g$,
\begin{equation}\label{eq:mumford-noether-relation}
\kMum \;=\; 12\lambda_1 \;-\; \delta
\qquad \textup{in } H^{2}(\overline{\cM}_g,\Q),
\end{equation}
in which $\kscal$ does \emph{not} appear.
\end{definition}

\begin{lemma}[Type discipline in the identity
$\mathrm{obs}_g=\kscal\cdot\lambda_g$; \ClaimStatusProvedHere]
\label{lem:theorem-D-type-discipline}
\index{Theorem D!type discipline|textbf}
In the identity
\[
\mathrm{obs}_g(\cA) \;=\; \kscal(\cA)\cdot\lambda_g
\qquad \textup{in } H^{2g}(\overline{\cM}_g,\Q),
\qquad \textup{(}UNIFORM-WEIGHT\textup{)},
\]
the left-hand side is a class in $H^{2g}$, the right-hand side is
the rational scalar $\kscal(\cA)$ acting on the top Hodge Chern class
$\lambda_g=c_g(\mathbb{E})\in H^{2g}$. Replacing $\kscal$ by $\kMum$
would produce the type-mismatched object $\kMum\cdot\lambda_g\in
H^{2+2g}$, which is a class of the \emph{wrong degree} and cannot
equal $\mathrm{obs}_g\in H^{2g}$.
\end{lemma}

\begin{proof}
Counting cohomological degrees: $\kscal(\cA)\in H^0=\Q$ and
$\lambda_g\in H^{2g}$, so $\kscal(\cA)\lambda_g\in H^{2g}$. By
contrast $\kMum\in H^{2}$ and $\kMum\lambda_g\in H^{2+2g}$, which is
the wrong degree. The identification
$\mathrm{obs}_g=\kscal\lambda_g$ is therefore internally consistent
only with the scalar reading of $\kappa$; the Mumford-class reading
is excluded by degree. \end{proof}

\begin{remark}[Style convention on $\kappa$ without decoration]
\label{rem:kappa-bare-convention}
\index{kappa@$\kappa$!bare-symbol convention}
Inside a displayed formula of the form $\kappa\cdot\lambda_g$ or
$\kappa\cdot S_g$ with $S_g\in H^{2g}$ a top-degree class, a bare
``$\kappa$'' always denotes the scalar $\kscal$ of
Definition~\ref{def:kappa-scalar-vs-mumford}(a): the type system
forces this reading. Mumford's class is denoted $\kMum$ or $\kappa_1$
throughout and never appears without its subscript. The
Miller--Morita--Mumford higher classes $\kappa_j$ for $j\geq 2$,
when they occur, are always subscripted. See
Appendix~\ref{app:notation-index} for the full table.
\end{remark}

\section{Genus ladder of the universal obstruction}
\label{sec:theta-A-genus-ladder}
\index{universal obstruction!genus ladder|textbf}
\index{Maurer--Cartan!genus-$g$ component|textbf}
\index{Theta_A@$\Theta_\cA^{(g)}$!ladder $g=1,2,3$|textbf}

We now exhibit the Maurer--Cartan element $\Theta_\cA$ of the
ordered convolution dGLA
\textup{(}Definition~\textup{\ref{def:ordered-convolution-dgla}}%
\textup{)} as a tower indexed by genus, together with the chain-level
witness for the induction step that controls the passage from genus
$g$ to genus $g+1$.

\subsection{Genus-$1$ element}
\label{subsec:theta-A-g1}

\begin{proposition}[Genus-$1$ MC element; \ClaimStatusProvedHere]
\label{prop:theta-A-genus1}
\index{Theta_A@$\Theta_\cA^{(1)}$!explicit elliptic form|textbf}
\index{Bernard--Felder propagator}
\index{Eisenstein series!genus-$1$ MC element}
On a marked elliptic curve $E_\tau=\C/(\Z+\tau\Z)$ with
$\tau\in\mathfrak{H}$, the genus-$1$ component
$\Theta_\cA^{(1)}\in\gAord$ of the universal MC element is
\begin{equation}\label{eq:theta-A-g1}
\Theta_\cA^{(1)}(\tau;z_{1},\ldots,z_{n})
\;=\;
\sum_{i<j} t_{ij}\,\eta^{(1)}_{\mathrm{BF}}(\tau, z_{ij})
\;\otimes\; \mathrm{pr}_{n}^{\mathrm{ord}},
\end{equation}
where $\eta^{(1)}_{\mathrm{BF}}(\tau,z)=\zeta_{\tau}(z)-z\,G_{2}(\tau)$
is the Bernard--Felder elliptic propagator \textup{(}\cite{Bernard88},
\cite{CEE10}\textup{)}, $t_{ij}\in\cA^{\otimes n}$ is the elementary
pairing tensor, and $\mathrm{pr}_{n}^{\mathrm{ord}}$ projects to the
ordered configuration space $\Confnord(E_\tau)$. The element
$\Theta_\cA^{(1)}$ satisfies the Maurer--Cartan equation
\begin{equation}\label{eq:mc-g1}
D\cdot\Theta_\cA^{(1)}
\;+\;
\tfrac{1}{2}\bigl[\Theta_\cA^{(1)},\Theta_\cA^{(1)}\bigr] \;=\; 0
\qquad \textup{in } \gAord\textup{[}1\textup{]},
\end{equation}
and its scalar curvature class is
\begin{equation}\label{eq:obs-g1}
[\mathrm{obs}_{1}(\cA)]
\;=\;
\kscal(\cA)\cdot\lambda_{1}
\qquad \textup{in } H^{2}(\overline{\cM}_{1,1},\Q).
\end{equation}
\end{proposition}

\begin{proof}
Equation~\eqref{eq:mc-g1} is the elliptic Arnold identity
$\sum_{\mathrm{cyc}} t_{ij}t_{jk}\,\eta^{(1)}_{\mathrm{BF}}(\tau,z_{ij})
\eta^{(1)}_{\mathrm{BF}}(\tau,z_{jk})=0$ modulo the exact term
$\bar\partial \eta^{(1)}_{\mathrm{BF}}=\pi\,\delta_{z=0}$, which is the
Calaque--Enriquez--Etingof proof of the KZB flatness
\textup{(}\cite[Theorem~4.1]{CEE10}\textup{)}.
Equation~\eqref{eq:obs-g1} is the genus-$1$ curvature extraction:
$\dfib^{\,2}|_{E_\tau}=\kscal(\cA)\,\omega_{E_\tau}$
and $H^{2}(\overline{\cM}_{1,1},\Q)=\Q\lambda_{1}$. \end{proof}

\subsection{Genus-$2$ element}
\label{subsec:theta-A-g2}

Three data organise the genus-$2$ MC element on $\overline{\cA_2}$:
the Beilinson--Drinfeld Green function $G_{\Sigma_2}(z,w)$ on the
smooth fibre, a Siegel Fourier coefficient $c_n(Z)$ at each admissible
Humbert divisor $H_n\subset\cA_2$ \textup{(}a Humbert divisor is the
locus where the polarised abelian surface $\mathrm{Jac}(\Sigma_2)$
acquires real multiplication of discriminant $n\equiv 0,1\pmod 4$;
\cite{vdG1988}, \cite{Humbert1899}\textup{)}, and the stress tensor
$T(z)$ of the conformal chiral algebra $\cA$. The first assembles a
bulk Arnold-relation two-form on $\Sigma_2\times\Sigma_2$; the second
supplies Mittag--Leffler counterterms on each Humbert stratum; the
third carries the conformal content of the bracket.

\begin{proposition}[Genus-$2$ MC element, explicit form;
\ClaimStatusProvedHere]
\label{prop:theta-A-genus2}
\index{Theta_A@$\Theta_\cA^{(2)}$!Beilinson--Drinfeld Green function|textbf}
\index{Theta_A@$\Theta_\cA^{(2)}$!Humbert Fourier expansion|textbf}
\index{Beilinson--Drinfeld Green function}
\index{Humbert divisor!Fourier coefficient $c_n$}
Fix a smooth genus-$2$ curve~$\Sigma_{2}$ with period matrix
$Z=\bigl(\begin{smallmatrix}\tau_1 & \tau_{12}\\ \tau_{12} & \tau_2
\end{smallmatrix}\bigr)\in\mathfrak{H}_2$ parametrising
$[\Sigma_{2}]\in\cM_{2}\hookrightarrow\cA_2\subset\overline{\cA_2}$.
Write $\omega_1,\omega_2\in H^{0}(\Sigma_2,\omega_{\Sigma_2})$ for
the normalised abelian differentials with
$\int_{\alpha_i}\omega_j=\delta_{ij}$, $\int_{\beta_i}\omega_j=Z_{ij}$,
and let $\theta(u;Z)$ be the Siegel theta function of characteristic
$[\,0\,;\,0\,]$ with $D\in\mathrm{Jac}(\Sigma_2)$ the Riemann vector.
The \emph{Beilinson--Drinfeld Green function}
\textup{(}\cite[§16.5]{FrenkelBenZvi}, \cite{Arakelov1974}\textup{)}
is the symmetric $(1,0)$-form in two variables
\begin{equation}\label{eq:bd-green-g2}
G_{\Sigma_2}(z,w)\;=\;
\partial_z\log\theta\!\bigl(\!\textstyle\int_{w}^{z}\omega-D;\,Z
\bigr)
\;-\;
2\pi i\sum_{k,l=1}^{2}(\mathrm{Im}\,Z)^{-1}_{kl}\,
\omega_k(z)\,\overline{\omega_l(w)},
\end{equation}
the unique meromorphic $(1,0)$-form satisfying
$\bar\partial_w G_{\Sigma_2}(z,w) = \pi\,\delta_{z=w}\,dz$ and
$\int_{\alpha_j}G_{\Sigma_2}(z,\cdot)=0$ for every $\alpha$-cycle.
Write the \emph{Humbert coefficient} at discriminant
$n\equiv 0,1\pmod 4$ as the Siegel Fourier expansion
\begin{equation}\label{eq:humbert-fourier}
c_n(Z)\;=\;
\sum_{\substack{(a,b,c)\in\Z_{\geq 0}^{3}\\ b^{2}-4ac=-n,\,\gcd(a,b,c)=1}}
q_1^{\,a}\,q_2^{\,c}\,r^{b},
\qquad
q_1=e^{2\pi i\tau_1},\ q_2=e^{2\pi i\tau_2},\ r=e^{2\pi i\tau_{12}},
\end{equation}
the Gritsenko lift of the primitive theta-series of the Humbert
quadratic form $b^2-4ac=-n$
\textup{(}\cite{Gritsenko1994}, \cite[Ch.~III]{vdG1988}\textup{)}.
Let $h_n\in\{1,2,3,4,\ldots\}$ index the admissible Humbert
divisors \textup{(}$n=1$: reducible Jacobian $E_1\times E_2$;
$n=4$: bielliptic; $n=5,8,9,\ldots$: real multiplication\textup{)}.
Then the genus-$2$ component of the universal MC element
\textup{(}Definition~\ref{def:ordered-convolution-dgla}\textup{)} is
\begin{equation}\label{eq:theta-A-g2}
\boxed{\;
\Theta_\cA^{(2)}(Z;z_{1},\ldots,z_{n})
\;=\;
\sum_{i<j} t_{ij}\,G_{\Sigma_2}(z_i,z_j)
\;+\;
\kscal(\cA)\sum_{n\geq 1}^{\mathrm{adm}}
c_n(Z)\cdot H_n\cdot T(z_{\star_n})
\;}
\end{equation}
where $t_{ij}\in\cA^{\otimes n}$ is the two-point Casimir, $H_n$
denotes the reduced Humbert divisor on $\cA_2$ \textup{(}pulled back
to $\cM_2$ via Torelli\textup{)}, and
$z_{\star_n}\in\Sigma_2$ is the base point selected by the $H_n$
period polarisation. Equation~\eqref{eq:theta-A-g2} satisfies the
Maurer--Cartan equation
\begin{equation}\label{eq:mc-g2}
D\cdot\Theta_\cA^{(2)}\;+\;\tfrac{1}{2}\bigl[\Theta_\cA^{(2)},
\Theta_\cA^{(2)}\bigr]\;=\;0
\qquad\textup{in }\gAord^{(2)}[1]
\end{equation}
on the whole of $\overline{\cA_2}$ after the Humbert counterterm of
equation~\eqref{eq:theta-A-g2} is included; the associated scalar
curvature class is
\begin{equation}\label{eq:obs-g2}
[\mathrm{obs}_{2}(\cA)]
\;=\;
\kscal(\cA)\cdot\lambda_{2}
\qquad \textup{in } H^{4}(\overline{\cM}_{2},\Q).
\end{equation}
\end{proposition}

\begin{proof}
\emph{Step 1 (bulk Arnold-relation bracket).} On
$\cM_2\setminus\bigcup_n H_n$ the form $G_{\Sigma_2}$ is smooth and
$\bar\partial_w G_{\Sigma_2}(z,w)=\pi\,\delta_{z=w}\,dz$. The
genus-$2$ analogue of Arnold's triple identity is the elliptic
sewing identity \cite[Prop.~2.4]{Enriquez2008}
\begin{equation}\label{eq:arnold-g2}
\sum_{\mathrm{cyc}(ijk)}
G_{\Sigma_2}(z_i,z_j)\wedge G_{\Sigma_2}(z_j,z_k)
\;=\;
\pi\bigl(\delta_{z_i=z_j}+\delta_{z_j=z_k}+\delta_{z_k=z_i}\bigr)
\cdot\omega_{\Sigma_2}(z_\bullet),
\end{equation}
where the right-hand side is the delta-distribution supported on the
pairwise diagonals, weighted by the canonical class
$\omega_{\Sigma_2}$ \textup{(}the
$-(2\pi i/g)\sum(\mathrm{Im}\,Z)^{-1}_{kl}\omega_k\bar\omega_l$
correction in \eqref{eq:bd-green-g2} absorbs the harmonic mode so
that \eqref{eq:arnold-g2} is pointwise and not only modulo harmonic
forms\textup{)}. Applying the bracket in the convolution dGLA
$\gAord^{(2)}$ and using $[t_{ij},t_{jk}]=t_{ijk}$ on triple
indices, the bulk contribution
$\sum_{i<j}t_{ij}G_{\Sigma_2}(z_i,z_j)$ gives
$\tfrac12[\,\cdot\,,\,\cdot\,]$ equal to the residue against the
diagonal $\omega_{\Sigma_2}$, which cancels exactly
$D\cdot\sum_{i<j}t_{ij}G_{\Sigma_2}(z_i,z_j)$ by the Stokes identity
$\bar\partial G_{\Sigma_2}=\pi\delta_{\mathrm{diag}}$.

\emph{Step 2 (Humbert counterterm).} At each admissible divisor
$H_n\subset\cA_2$ the abelian surface develops real multiplication
by $\mathcal{O}_{\Q(\sqrt{n})}$; the period coordinate
$\tau_{12}/\sqrt{\tau_1\tau_2-\tau_{12}^2}$ satisfies a rational
relation with discriminant $n$, and
$G_{\Sigma_2}$ acquires an extra polar term in the transverse
Siegel coordinate $\xi_n$ with residue
$\mathrm{Res}_{H_n}G_{\Sigma_2}=c_n(Z)\cdot\omega_{\Sigma_2}$
\textup{(}the Fourier expansion \eqref{eq:humbert-fourier} is the
Gritsenko lift of the primitive theta-series of the binary
quadratic form $b^2-4ac=-n$; the coefficient at $q_1^a q_2^c r^b$
counts primitive representations\textup{)}. The scalar coefficient
$\kscal(\cA)$ appears because the conformal vector $T$ carries
central charge $c=2\kscal(\cA)$ in the bracket $[T,T]$ \textup{(}see
\cite[§5]{KRW88-convention} for the Sugawara normalisation\textup{)}.
Hence the Humbert counterterm $\kscal(\cA)\sum_n c_n(Z)\,H_n\,T$
cancels the transverse-Siegel residue of the bulk bracket on each
$H_n$, which is the Mittag--Leffler tier-$1$ input of
Remark~\ref{rem:theta-A-humbert-tier1}.

\emph{Step 3 (scalar curvature).} The bar-Hodge metric on
$\det\pi_*B^{\mathrm{ord}}(\cA)$ coincides \textup{(}Bismut--Gillet--Soul\'e,
\cite[Thm.~0.1]{BGS1988}\textup{)} with the Quillen metric on
$\det\mathbb{E}$ up to a scalar factor $\kscal(\cA)$
\textup{(}Proposition~\ref{prop:scalar-obstruction-hodge-euler}\textup{)};
the first Chern class of $\det\mathbb{E}$ on $\overline{\cM}_2$ is
$\lambda_1$, and its top-degree image in $H^4$ is $\lambda_2$. Hence
$[\mathrm{obs}_2(\cA)]=\kscal(\cA)\lambda_2$.
\end{proof}

\begin{corollary}[Three independent verifications of
$\Theta_\cA^{(2)}$; \ClaimStatusProvedHere]
\label{cor:theta-A-g2-three-paths}
\index{verification!three paths for $\Theta_\cA^{(2)}$|textbf}
The Maurer--Cartan identity \eqref{eq:mc-g2} is verified along three
mutually independent paths.

\textup{(i) \emph{Direct Arnold bracket}.} Inserting
\eqref{eq:bd-green-g2} into \eqref{eq:arnold-g2} and contracting
against $t_{ij}t_{jk}=t_{ijk}\in\cA^{\otimes 3}$ reduces the triple
product to a sum of Kronecker delta-distributions on pairwise
diagonals; combined with the Humbert residue
$\mathrm{Res}_{H_n}G_{\Sigma_2}=c_n(Z)\omega_{\Sigma_2}$ this cancels
$D\Theta_\cA^{(2)}$ term by term. No cohomological input.

\textup{(ii) \emph{Kontsevich formality degeneration}
$\mathrm{Vol}(E_2)\to 0$}. Restrict $Z$ to the Humbert divisor $H_1$
\textup{(}reducible Jacobian $E_{\tau_1}\times E_{\tau_2}$,
$\tau_{12}=0$\textup{)}; the Green function factorises as
$G_{\Sigma_2}\vert_{H_1}=G_{E_{\tau_1}}\oplus G_{E_{\tau_2}}$ with
$G_{E_\tau}=\eta^{(1)}_{\mathrm{BF}}(\tau,z)$ the Bernard--Felder
propagator of Proposition~\ref{prop:theta-A-genus1}. Kontsevich
formality on $E_{\tau_1}\times E_{\tau_2}$
\textup{(}\cite[§6]{CEE10}\textup{)} identifies
$\Theta_\cA^{(2)}\vert_{H_1}=\Theta_\cA^{(1)}(\tau_1)\otimes
\mathrm{id}+\mathrm{id}\otimes\Theta_\cA^{(1)}(\tau_2)$ in the
limit $\mathrm{Vol}(E_2)=\mathrm{Im}\,\tau_2\to 0$; the MC equation
at $g=2$ pulls back to the sum of the two $g=1$ MC equations, each
of which is \eqref{eq:mc-g1}. This is independent of path~(i) since
it uses no direct triple-bracket computation.

\textup{(iii) \emph{Mumford $\kappa$-class specialisation}}. The
socle integral $\int_{\overline{\cM}_2}\lambda_2\lambda_1=1/5760$
\textup{(}\cite{Mumford83}\textup{)} combined with the
Grothendieck--Riemann--Roch identity
$\lambda_2=\lambda_1\kappa_1/12+\text{boundary}$ on
$\overline{\cM}_2$ gives
\begin{equation}\label{eq:obs-g2-archetype-check}
\int_{\overline{\cM}_2}[\mathrm{obs}_2(\cA)]\,\lambda_1
\;=\;
\frac{\kscal(\cA)}{5760},
\end{equation}
matching Proposition~\ref{prop:archetype-specialisation-obs-g}
evaluations at the four archetypes:
$k/5760$ \textup{(}Heisenberg\textup{)},
$\dim(\fg)(k+\hv)/(2\hv\cdot 5760)$ \textup{(}affine KM\textup{)},
$c/11520$ \textup{(}Virasoro\textup{)},
$(6\lambda^2-6\lambda+1)/5760$ \textup{(}$\beta\gamma_\lambda$\textup{)}.
This uses no bar-level bracket and is purely a cohomological
Hodge-integral check. \qed
\end{corollary}

\subsection{Genus-$3$ element and beyond}
\label{subsec:theta-A-g3}

\begin{proposition}[Genus-$3$ MC element; \ClaimStatusProvedHere]
\label{prop:theta-A-genus3}
\index{Theta_A@$\Theta_\cA^{(3)}$!Schottky locus|textbf}
\index{Schottky locus!genus-$3$ MC element}
On a non-hyperelliptic genus-$3$ curve $\Sigma_{3}$, the
Beilinson--Drinfeld Green function $G_{\Sigma_{3}}$ decomposes under
the Torelli map $\mathrm{Tor}\colon\cM_{3}\to\cA_{3}$ as
$G_{\Sigma_{3}}=G_{\mathrm{Jac}}|_{\mathrm{Tor}(\Sigma_{3})}+
R_{\Schottky}$, where $R_{\Schottky}$ is the Schottky correction
supported on $\cM_{3}\setminus\mathrm{Tor}^{-1}(\Schottky\text{-locus})$.
The genus-$3$ MC element
\begin{equation}\label{eq:theta-A-g3}
\Theta_\cA^{(3)}
\;=\;
\sum_{i<j} t_{ij}\bigl(G_{\mathrm{Jac}}+R_{\Schottky}\bigr)(z_{i},z_{j})
\;+\;
\sum_{i}\bigl(H_{i}^{(2)}+H_{i}^{\Schottky}\bigr)T(z_{i})
\end{equation}
satisfies
$D\Theta_\cA^{(3)}+\tfrac{1}{2}[\Theta_\cA^{(3)},\Theta_\cA^{(3)}]=0$,
and
\begin{equation}\label{eq:obs-g3}
[\mathrm{obs}_{3}(\cA)]
\;=\;
\kscal(\cA)\cdot\lambda_{3}
\qquad \textup{in } H^{6}(\overline{\cM}_{3},\Q).
\end{equation}
\end{proposition}

\begin{proof}
The Schottky decomposition is Schottky's identity
$G_{\Sigma_{3}}-\mathrm{Tor}^{*}G_{\mathrm{Jac}}=R_{\Schottky}$
\textup{(}explicit on the theta-divisor of
$\mathrm{Jac}(\Sigma_{3})$\textup{)}. The MC equation reduces to the
genus-$2$ MC equation on the boundary
$\partial\overline{\cM}_{3}\supset\xi_{h}(\overline{\cM}_{2,1}
\times\overline{\cM}_{1,1})\cup\xi_{\mathrm{irr}}(\overline{\cM}_{2,2})$
via the clutching-compatibility
Theorem~\ref{thm:separating-nonseparating-clutching-compatibility}.
Equation~\eqref{eq:obs-g3} is
Proposition~\ref{prop:scalar-obstruction-hodge-euler} at $g=3$.
\end{proof}

\subsection{Induction step: $g \Rightarrow g+1$}
\label{subsec:theta-A-induction}

\begin{theorem}[Induction step for the universal MC ladder;
\ClaimStatusProvedHere]
\label{thm:theta-A-induction}
\index{induction!genus $g\Rightarrow g+1$ for MC element|textbf}
\index{Mittag--Leffler!tower for MC element|textbf}
\index{MC element!induction step with explicit witness|textbf}
Assume the inductive data:
\begin{itemize}
\item $\Theta_\cA^{(h)}\in\gAord^{(h)}$ satisfying
$D\Theta_\cA^{(h)}+\tfrac{1}{2}[\Theta_\cA^{(h)},\Theta_\cA^{(h)}]=0$
for all $0\leq h\leq g$;
\item compatibility with both clutching maps:
$\xi_{h}^{*}\Theta_\cA^{(g)}=\Theta_\cA^{(h)}\boxtimes
\Theta_\cA^{(g-h)}$ and $\xi_{\mathrm{irr}}^{*}\Theta_\cA^{(g)}=
\Delta_{\mathrm{ns}}(\Theta_\cA^{(g-1)})$.
\end{itemize}
Then there exists $\Theta_\cA^{(g+1)}\in\gAord^{(g+1)}$ satisfying
$D\Theta_\cA^{(g+1)}+\tfrac{1}{2}[\Theta_\cA^{(g+1)},
\Theta_\cA^{(g+1)}]=0$, uniquely up to gauge, such that
\begin{enumerate}[label=\textup{(\roman*)}]
\item $\xi_{h}^{*}\Theta_\cA^{(g+1)}
=\Theta_\cA^{(h)}\boxtimes\Theta_\cA^{(g+1-h)}$ for every
$1\leq h\leq\lfloor(g+1)/2\rfloor$;
\item $\xi_{\mathrm{irr}}^{*}\Theta_\cA^{(g+1)}=
\Delta_{\mathrm{ns}}(\Theta_\cA^{(g)})$;
\item the scalar curvature class is
$[\mathrm{obs}_{g+1}(\cA)]=\kscal(\cA)\cdot\lambda_{g+1}$
in $H^{2g+2}(\overline{\cM}_{g+1},\Q)$.
\end{enumerate}
The compatibility witness is the Mittag--Leffler tower
$\{\gAord^{(g,N)}\}_{N\geq 1}$ of weight-truncated convolution
dGLAs, in which the truncation differential
$d_{\mathrm{ML}}^{(N)}\colon\gAord^{(g+1,N)}\to\gAord^{(g+1,N-1)}$
is compatible with the boundary restrictions $\xi_{h}^{*}$ and
$\xi_{\mathrm{irr}}^{*}$ by construction.
\end{theorem}

\begin{proof}
The obstruction to solving the MC equation at genus $g+1$ lies in
$H^{1}(\gAord^{(g+1)})$, and its boundary-restriction vanishes on
both clutching images by the inductive hypothesis. By the
Graber--Vakil vanishing
theorem~\cite{GraVak05} (Step~3 of the proof of
Theorem~\ref{thm:clutching-uniqueness-socle-projection}), an element
of $H^{1}(\gAord^{(g+1)})$ vanishing on every boundary divisor lies
in the socle of the tautological ring. Pairing against a fixed
admissible socle class $\omega_{\mathrm{soc}}$ gives the rational
number $\pi_{\mathrm{soc}}(\mathrm{obs}_{g+1}/\kscal(\cA))$; by the
genus-$1$ normalisation and the uniform-weight
Proposition~\ref{prop:scalar-obstruction-hodge-euler}, this socle
projection equals $\pi_{\mathrm{soc}}(\lambda_{g+1})$. Hence the
obstruction is socle-trivial and the Mittag--Leffler tower
$\{\gAord^{(g+1,N)}\}$ converges
\textup{(}$\mathrm{lim}^{1}\Theta_\cA^{(g+1,N)}=0$ is the
Mittag--Leffler condition verified in
Proposition~\textup{\ref{prop:mittag-leffler-tower-bar-convergence}}%
\textup{)} to a solution $\Theta_\cA^{(g+1)}$ satisfying (i)--(iii).
Uniqueness up to gauge is the standard first-order deformation-theory
argument for curved MC elements~\cite{Kontsevich97,Costello11}.
\end{proof}

\begin{remark}[Tier-$1$ connection to Humbert neighbourhoods]
\label{rem:theta-A-humbert-tier1}
\index{Humbert!neighbourhoods in induction step}
The Mittag--Leffler truncation $d_{\mathrm{ML}}^{(N)}$ is
exactly the finite-weight approximation to the obstruction
integrand in a neighbourhood of the Humbert locus
$\Humbert\subset\cM_{2}$, where the non-Jacobian/non-Schottky
corrections concentrate. At genus $2$ this is $H_{i}T(z_{i})$ of
equation~\eqref{eq:theta-A-g2}; at higher genus it is the recursive
lift of the Humbert correction across the clutching maps. The
compatibility witness of Theorem~\ref{thm:theta-A-induction}
therefore ties directly into the Tier-$1$ programme of
Mittag--Leffler analysis across Humbert strata.
\end{remark}

\section{Three verification paths for Theorem~D}
\label{sec:theorem-D-three-paths}
\index{Theorem D!three independent verification paths|textbf}
\index{verification!multi-path for $\mathrm{obs}_g=\kappa\lambda_g$}

The identity $\mathrm{obs}_g(\cA)=\kscal(\cA)\cdot\lambda_g$ is
verified along three mutually independent paths; we record the
verification explicitly at the leading genera.

\subsection{Path (i): direct MC check at $g=1$}

\begin{proposition}[$g=1$ direct MC verification;
\ClaimStatusProvedHere]
\label{prop:mc-direct-g1-verification}
\index{verification!direct MC at genus 1}
For every archetype $\cA\in\{\cH_{k},V_{k}(\fg),\beta\gamma_\lambda,
\Vir_{c}\}$, equation~\eqref{eq:mc-g1} is verified by the explicit
elliptic bracket identity
\begin{equation}\label{eq:mc-verification-direct}
\sum_{\mathrm{cyclic}(ijk)}
\eta^{(1)}_{\mathrm{BF}}(\tau,z_{ij})\,
\eta^{(1)}_{\mathrm{BF}}(\tau,z_{jk})
\;=\;
0
\qquad \textup{modulo } \pi\,\delta_{\mathrm{diag}},
\end{equation}
which is the elliptic Arnold identity of Calaque--Enriquez--Etingof
\textup{(}\cite[Section~4]{CEE10}\textup{)}. The scalar coefficient
extracted from the curvature is
$\kscal(\cH_{k})=k$,
$\kscal(V_{k}(\fg))=\dim(\fg)(k+\hv)/(2\hv)$,
$\kscal(\Vir_{c})=c/2$,
$\kscal(\beta\gamma_\lambda)=6\lambda^{2}-6\lambda+1$.
The Hodge class is $\lambda_{1}=\delta/12\in H^{2}(\overline{\cM}_{1,1},\Q)$
by Mumford's Noether formula \textup{(}\cite{Mumford83}\textup{)};
both sides of $\mathrm{obs}_{1}=\kscal\lambda_{1}$ therefore evaluate
to $\kscal(\cA)/24$ against the generator
$[\overline{\cM}_{1,1}]\in H_{2}(\overline{\cM}_{1,1},\Q)$.
\end{proposition}

\subsection{Path (ii): Grothendieck--Riemann--Roch}

\begin{proposition}[GRR verification at all $g$;
\ClaimStatusProvedHere]
\label{prop:grr-verification-all-g}
\index{verification!Grothendieck--Riemann--Roch}
Mumford's Grothendieck--Riemann--Roch~\cite{Mumford83} applied to
$\pi\colon\overline{\cC}_{g}\to\overline{\cM}_{g}$ with the relative
dualising sheaf $\omega_{\pi}$ gives the Chern character identity
\begin{equation}\label{eq:grr-mumford}
\mathrm{ch}(\pi_{*}\omega_{\pi})
\;=\;
\pi_{*}\bigl(\mathrm{ch}(\omega_{\pi})\cdot\mathrm{Td}(T_{\pi})\bigr)
\;=\;
\sum_{j\geq 0}\frac{B_{2j}}{(2j)!}\,\kappa_{2j-1}
\qquad \textup{in } H^{*}(\overline{\cM}_{g},\Q),
\end{equation}
where $\kappa_{j}=\pi_{*}c_{1}(\omega_{\pi})^{j+1}$ are the
Miller--Morita--Mumford classes \textup{(}Mumford~$1983$,
Miller~$1986$, Morita~$1987$\textup{)}. The
degree-$g$ component of $\mathrm{ch}(\mathbb{E})$ is
$(-1)^{g-1}\mathrm{ch}_{g}(\mathbb{E})$, and summing with the
signs of the lambda class produces
$\lambda_{g}=c_{g}(\mathbb{E})$ after Newton's identities. The
bar-curvature class $\mathrm{obs}_{g}$ coincides with
$\kscal(\cA)\cdot c_{g}(\mathbb{E})$ by the Bismut--Gillet--Soul\'e
family-index identification of the bar-Hodge metric with the
Quillen metric on $\mathbb{E}$
\textup{(}Proposition~\textup{\ref{prop:scalar-obstruction-hodge-euler}}%
\textup{)}. The appearance of $\kappa_{j}$ in
equation~\eqref{eq:grr-mumford} is the Mumford-class (not chiral
$\kscal$) by Remark~\ref{rem:kappa-bare-convention}.
\end{proposition}

\subsection{Path (iii): specialisation to the four archetypes}

\begin{proposition}[Archetype specialisation;
\ClaimStatusProvedHere]
\label{prop:archetype-specialisation-obs-g}
\index{verification!archetype specialisation Heisenberg/Virasoro/KM}
For each archetype, the identity
$\mathrm{obs}_g(\cA)=\kscal(\cA)\lambda_g$ evaluates against the
socle class $\omega_{\mathrm{soc}}=\lambda_{g-1}\lambda_{g-2}$ to
the Faber--Pandharipande Hodge integrals
$\int_{\overline{\cM}_g}\lambda_g\lambda_{g-1}\lambda_{g-2}$
\textup{(}\cite{FP00}\textup{)}. At $g=2$ with
$\omega_{\mathrm{soc}}=\lambda_{1}$
\textup{(}Corollary~\ref{cor:genus-2-explicit-match}\textup{)}:
\begin{align*}
\int_{\overline{\cM}_2}\mathrm{obs}_{2}(\cH_{k})\lambda_{1}
&= k/5760,\\
\int_{\overline{\cM}_2}\mathrm{obs}_{2}(V_{k}(\fg))\lambda_{1}
&= \dim(\fg)(k+\hv)/(2\hv\cdot 5760),\\
\int_{\overline{\cM}_2}\mathrm{obs}_{2}(\Vir_{c})\lambda_{1}
&= c/11520.
\end{align*}
At $g=3$ the analogous evaluation against
$\lambda_{2}\lambda_{1}$ gives
$\kscal(\cA)/2903040$ via Faber's genus-$3$ Hodge integral
$\int_{\overline{\cM}_3}\lambda_3\lambda_2\lambda_1=1/2903040$
\textup{(}\cite{FP00}\textup{)}.
Each specialisation is an independent numerical check on the
scalar-lane Theorem~D.
\end{proposition}

\begin{proof}
At $g=2$ the socle integral
$\int_{\overline{\cM}_2}\lambda_{2}\lambda_{1}=1/5760$ is
Mumford's~\cite{Mumford83} classical computation
(proved via $12\lambda_{1}=\delta$ on
$\overline{\cM}_{1,1}$, Mumford's relation
$c(\mathbb{E})c(\mathbb{E}^\vee)=1$ which gives
$\lambda_{2}=\lambda_{1}\kappa_{1}/12$ modulo boundary on
$\overline{\cM}_{2}$, and direct evaluation).
Substitution of the archetype $\kscal$ values from
Theorem~\ref{thm:genus-universality}(i) gives the stated numerical
outcomes. At $g=3$ the Hodge integral
$\int_{\overline{\cM}_3}\lambda_3\lambda_2\lambda_1=1/2903040$ is
from~\cite{FP00}; the Virasoro
case specialises to $c/5806080$. \end{proof}

\begin{remark}[Three paths are genuinely independent]
\label{rem:three-paths-independence}
\index{independence!three verification paths for Theorem D}
Path (i) is chain-level, local to an elliptic fiber, and uses only
the Bernard--Felder propagator and the Arnold identity. Path (ii) is
cohomological, global on $\overline{\cM}_g$, and uses
Grothendieck--Riemann--Roch and the Bismut--Gillet--Soul\'e anomaly
formula. Path (iii) is numerical, evaluates the socle pairing, and
uses the Faber--Pandharipande Hodge integrals. No theorem or
primary lemma is common to any two paths; the cited inputs are
disjoint. This realises the three-independent-paths discipline
established in the programme constitution.
\end{remark}
